% ====================================================================
\section{Formal verification in Lean}\label{sec:lean}
% ====================================================================
\begin{sloppypar}

The complete argument of this paper has been formalized in Lean~4 (toolchain \texttt{v4.33.1}) on Mathlib revision
\texttt{0df444a3}, the same pins as the Google DeepMind \emph{formal-conjectures} repository at commit
\texttt{c77fdd446d23539302aa73bcfae4a9fa962b8382}. The formalization is about $9{,}800$ lines.

\paragraph{The target statement.} The theorem proved is \texttt{Erdos1088.erdos\_1088}:
\[
\forall n\ge 3,\qquad \bigl(d\mapsto \log_2 f(d,n)\bigr) = o(d)\quad (d\to\infty),
\]
where \texttt{f d n} is defined as $\inf\{m\ge n:$ every $m$-point subset of $\R^d$ contains an $n$-point subset
with $\binom n2$ distinct distances$\}$. The statement and all definitions it uses (\texttt{f}, \texttt{cardSet},
\texttt{HasSubsetWithDistinctDistances}, \texttt{PairwiseDistancesDistinct}, \texttt{distinctDistances},
\texttt{distanceSet} and the notation $\R^d$) are copied character for character from formal-conjectures. A script
checks this against the upstream files. Only module headers and metadata attributes are dropped.

Lean's \texttt{sInf} of an empty set of naturals is $0$, so a proof could in principle exploit an empty
\texttt{cardSet}. The formal proof does not do this: it exhibits an explicit element of \texttt{cardSet d n} and
proves $f(d,n)\le n\cdot 2^{C_n\sqrt d\log_2(d+2)}$, which is the Main Theorem.

\paragraph{Trust base.} \texttt{\#print axioms Erdos1088.erdos\_1088} reports only \texttt{propext},
\texttt{Classical.choice} and \texttt{Quot.sound}, the standard axioms of Lean and Mathlib. The development
contains no \texttt{sorry}, no \texttt{axiom} declarations and no \texttt{native\_decide}.

\paragraph{Differences from the informal proof.} The formalization follows Sections~\ref{sec:LR}--\ref{sec:main}
with the following changes, none of which affects the Main Theorem.
\begin{itemize}[leftmargin=1.5em]
\item The normalization step (NF3) of Lemma~\ref{lem:NF} is replaced by a tournament argument. As a result, the formal
  version of Theorem~\ref{thm:LR} has loss $2^{O(k)}\sqrt r\log r$ instead of $O(k^3\sqrt r\log r)$; for fixed $k$ the
  shape $2^{-O_k(\sqrt r\log r)}$ is the same, and this is all the later sections use.
\item John's theorem in Lemma~\ref{lem:R} is replaced by a factorization derived from the Kiefer--Wolfowitz
  optimal-design theorem, proved from the matrix determinant lemma and the spectral theorem. Sheppard's formula is
  proved from the rotation invariance of the standard Gaussian.
\item Several constants are larger. These include those in the Core Lemma and in Theorem~\ref{thm:Rplus}, the
  number of big radii at which Step~A applies ($3$ instead of $2$), and $\kst$. The case $n=3$ is derived from G1
  rather than from Blokhuis's theorem.
\item Lean made explicit some hypotheses that the text leaves implicit, for example $\sigma>0$ and $\Gamma\ge 0$ in
  the kill step, and $B\subseteq X$ in the Core Lemma. It found no gap in the argument.
\end{itemize}

\paragraph{Availability.} The Lean source, build instructions and audit scripts are available at
\url{https://github.com/H0W4RD-dev/erdos-1088-lean}.
\end{sloppypar}
